\section{Mon (Monoides Finitos)}

Os monoides representam a estrutura básica de composição e computação sequencial. Um monoide combina um conjunto com uma operação binária que deve ser associativa e possuir um elemento neutro (identidade).

\subsection{Objetos}
A modelação computacional de monoides finitos é realizada através de Tabelas de Cayley. Esta matriz armazena todos os resultados possíveis da operação binária entre os elementos do conjunto.

\begin{lstlisting}[caption={Tabela de Cayley para Monoide C2}]
% Elementos: 1 (Identidade), 2
Tabela = [1 2; 
          2 1];
e = 1; % Índice do elemento neutro
\end{lstlisting}

\subsection{Morfismos}
Os morfismos de monoides (homomorfismos) devem preservar a estrutura operacional: a imagem do produto de dois elementos deve ser igual ao produto das suas imagens, e a identidade deve ser mapeada na identidade.

\begin{lstlisting}[caption={Validação de Homomorfismo}]
f = [1, 2];
% 1. Verificar se f(e_A) == e_B
% 2. Verificar se f(TabelaA(i,j)) == TabelaB(f(i), f(j))
\end{lstlisting}